\section{A controlled learner with discovered quadratic directions}
\label{sec:experiments}

We next ask whether the distinction between raw target moments and residual
moments improves an actual trained learner. This study implements the
mechanisms as conventional ingredients; it is not a new spectral algorithm.
No candidate receives a teacher direction. The principal negative result is
that a better initial direction need not retain its advantage after fully
trainable refitting, and an apparent final-risk benefit is sensitive to a
simple replay control.

\subsection{Frozen populations, data access and optimization}
All inputs have distribution $N(0,I_{12})$. Two quadratic conditional means
are $q_A(x)=x^\top A x-\operatorname{tr}(A)$, with nonzero eigenvalues
$(.8,.45)$ and $(.65,-.5,.3,-.2)$, respectively. The first has no observation
noise; the second has independent Gaussian noise with standard deviation
$.3$. A null task has mean zero and noise standard deviation $.5$.
The nonlinear task uses a hidden orthogonal rotation $z=U^\top x$ and mean
$.8\tanh(1.5z_1)\tanh(1.5z_2)+.4\sin z_3$, with noise standard deviation
$.2$. Task rotations are fixed independently of run seeds and hidden from
all learners. All labels, teacher parameters and evaluator inputs are saved.

The learned model is
\[
 f_{W,a,c}(x)=c+\sum_{j=1}^k a_j\{(w_j^\top x)^2-\|w_j\|^2\}.
\]
Every hidden vector, output coefficient and intercept trains. Initial
vectors are normalized Gaussian draws and initial coefficients have standard
deviation $.01$. Seven arrivals each reveal 1,024 labels. The current block
and an old reservoir form the fitting pool; the block enters the reservoir
only after fitting. Thus all 7,168 labels arrive at the same times, with no
future data granted to the fixed comparator. Growing arms start at width
one and receive six scheduled additions, ending at seven. This design isolates
discovered directions; it does not claim autonomous width selection or safe
rejection. New coefficients start at zero, preserving the current prediction
before ordinary, potentially risk-increasing training.

Seven arms compare: historical-moment eigenbirth, old-reservoir residual
 eigenbirth, fresh-block residual eigenbirth, four learned residual probes,
random growth with reservoirs of 128 and 138 examples, and a fixed width-seven
online learner with 138-example replay. The historical arm stores 128 examples
and $S_n=n^{-1}\sum_i y_i(x_ix_i^\top-I)/2$, and selects the eigenvector of
largest absolute eigenvalue of $S_n-B$, where $B=\sum_j a_jw_jw_j^\top$.
Residual eigenbirth instead forms
$\{\overline{(y-f)x x^\top}-\overline{(y-f)}I\}/2$ on its permitted
old reservoir or current block. It therefore includes the current intercept
term. The zero-intercept theoretical variance formulas are not claimed as
exact variance identities for every fitted checkpoint.

Learned probes maximize squared empirical residual covariance divided by
candidate second moment plus $.001$, with twenty steps and four random
starts. This is an established residual-feature objective. Directions are
selected from observations and then all network parameters are refitted.
Random additions use the same type of ordinary initial direction. The fixed
learner knows the prescribed capacity cap, not the teacher rank.

Each arrival permits two million arithmetic-proxy units, including moment
accumulation, candidate passes, eigendecomposition and fitting. Saved search
work becomes additional training. Full-gradient work is charged at
$3dk+4k+3$ per example; matrix and probe calculations have separate explicit
charges in the runner. This is not matched hardware FLOPs. Persistent Adam
uses learning rate $.015$, minibatches of at most 128, coefficients
$(.9,.999)$, denominator constant $10^{-8}$, and gradient clipping at norm
20. New-unit optimizer moments are zero; existing optimizer state persists.
No hyperparameter was changed using development evaluator outcomes.

Four development seeds on two tasks checked execution. The source,
confirmation configuration and protocol were frozen at 05:40:48 UTC on
20 September 2026, before seeds 820000--820011. The resulting 48 paired
configurations comprise 336 trajectories and 2,352 primary events, completed
in 21.47 seconds using two CPU workers. Quadratic-task risks are evaluated
exactly as $2\|B-A\|_F^2+c^2+\sigma^2$; nonlinear-task risk uses 32,768
independent evaluator inputs per seed and integrates label noise. Every
trajectory is finalized before its evaluation. All reported intervals are
descriptive 95\% Student-$t$ intervals across twelve independent seeds,
without multiplicity adjustment.

\subsection{Final risks and a failed attribution}
\begin{table}[tb]
\centering\small
\begin{tabular}{lrrrr}
\toprule
Method & Null & Rank two & Indefinite & Nonlinear \\
\midrule
Historical moment & 0.260725 & 0.000017 & 0.096987 & 0.161909 \\
Replay residual & 0.258093 & 0.000013 & 0.100358 & 0.164914 \\
Fresh residual & 0.258299 & 0.000044 & 0.102353 & 0.165770 \\
Learned probe & 0.256781 & 0.000087 & 0.106314 & 0.166361 \\
Random R128 & 0.258551 & 0.000017 & 0.097090 & 0.160259 \\
Random R138 & 0.257820 & 0.000021 & 0.102128 & 0.162947 \\
Fixed width 7 & 0.256951 & 0.000014 & 0.116278 & 0.162604 \\
\bottomrule
\end{tabular}

\caption{Mean final population MSE, generated from the frozen raw results.
All methods finish at width seven. The rank-two task is noiseless; the other
quadratic noise floors are $.25$ and $.09$. Nonlinear risks here use the
prespecified Monte Carlo evaluator. Full intervals and every paired contrast
are supplied as CSV.}
\label{tab:quadratic-results}
\end{table}

On the indefinite task, historical moments appear favorable against fresh
residuals: their paired final-risk difference is $-.005366$, with interval
$[-.008879,-.001853]$. Their difference from the 138-example random-growth
arm is $-.005141\;[-.008227,-.002055]$. Those comparisons alone would support
an overly strong interpretation. The 128-example random-growth arm also
beats the 138-example arm by $-.005038\;[-.007834,-.002243]$, while its mean
risk is $.097090$, almost the same as historical moments at $.096987$.
Their exploratory paired historical-minus-random-128 difference is
$-.000103\;[-.001961,.001756]$.
The near memory-matched contrast therefore does not isolate useful historical
information. Independent code review found a concrete batching confound:
the 128-example reservoir creates a 1,152-row fitting pool, divisible by the
128-row minibatch size; the 138-example reservoir creates 1,162 rows and an
additional ten-row tail batch every epoch. The two random arms therefore
perform 975 versus 1,034 Adam updates despite the same 124,285 fit-example
accesses. Changes in update count, gradient noise and replay sample accompany
the storage change. These results must not be interpreted as evidence that
more memory harms learning. A separately frozen full-minibatch follow-up
below tests this identified confound.

On the noiseless rank-two task, every method reaches a small mean error,
ranging from $1.25\times10^{-5}$ to $8.73\times10^{-5}$. The historical
comparison with random 138 has interval
$[-2.335,1.608]\times10^{-5}$; no final advantage is resolved. On the null
task, all methods retain excess error above $.25$, illustrating that the
scheduled-width study does not itself prevent needless capacity or
optimization error. On the nonlinear task, historical moments versus random
138 has difference $-.001038\;[-.006336,.004260]$, again unresolved.
Intervals containing zero are not equivalence statements.

The fixed width-seven learner has competitive mean risk on three tasks and
a worse, variable mean on the indefinite task. The historical-minus-fixed
indefinite interval, $[-.039426,.000845]$, includes zero. This is a competent
online baseline under the specified optimizer and budget, not proof that
fixed models are optimal or that growing models uniformly dominate them.

\begin{figure}[tb]
\centering\includegraphics[width=\linewidth]{quadratic_trajectories.pdf}
\caption{Complete trained histories from saved outcomes. Bands are one
standard error across independent runs. All growing arms have the same
scheduled widths; improvements here concern direction and fitting mechanisms.
The rank-two panel uses a logarithmic risk axis.}
\end{figure}

\subsection{Same-checkpoint directions and finite-time refitting}
Lifetime histories diverge, so their final comparison cannot diagnose the
quality of identical candidate sets. We therefore froze a separate,
evaluator-only diagnostic. At every pre-growth checkpoint of the random-138
trajectory, all five direction mechanisms see the same current predictor,
optimizer state, retained rows and available data prefix. Each direction is
then followed by identical fitting work and minibatch randomness. Candidate
work is additional and logged, rather than claimed equal. Reconstructing a
historical moment from the archive explicitly charges all prefix rereads;
it does not purchase those labels again or retroactively give that storage
to the primary random learner. This adds 1,440 trained diagnostic branches.

For a quadratic task, the best possible output coefficient along a unit
vector $w$ yields gain $2\{w^\top(A-B)w\}^2$. This is computed from teacher
parameters only after training and never enters selection. Averaged first
within each seed over six checkpoints, historical directions exceed random
directions in this potential gain by $.08041\;[.04305,.11777]$ on rank-two
data and $.13116\;[.11530,.14702]$ on indefinite data. These are strong
initial direction differences.

After identical fully trainable refitting, the corresponding
historical-minus-random risk differences are
$2.713\times10^{-5}$, with interval $[-1.111,6.536]\times10^{-5}$, and
$-.002687\;[-.006656,.001282]$, respectively.
Neither interval resolves a refitted advantage. Better direction alignment
therefore does not by itself establish better learned models at the available
training budget. This directly limits the proposed mechanism, rather than
hiding a negative result behind a successful population direction formula.
The trained hidden vectors may move substantially; this diagnostic does not
claim that every source of erasure has been isolated.

\begin{figure}[tb]
\centering\includegraphics[width=\linewidth]{quadratic_checkpoint.pdf}
\caption{Same-checkpoint diagnostics. Top: evaluator-only best gain along
the supplied-to-refitting, data-discovered direction. Bottom: actual risk after
identical nonlinear refitting. These are different objects, and the large
initial direction gap need not remain a final learning advantage.}
\end{figure}

\subsection{Memory, computation and scientific limits}
Every primary arm spends approximately fourteen million arithmetic-proxy
units over its lifetime. Raw label accesses for fitting, proposals and moments
are saved by sample ID. Historical retention uses
$128(12+2)+144=1936$ eight-byte-word equivalents, including example labels and
identities; random 138 uses 1932. At final width seven, both additionally use
276 model/Adam words and six controller-state words. This comparison is nearly
matched in persistent algorithmic storage, not peak process memory. Current
blocks, duplicated fitting pools, candidate matrices, hidden optimizer state
and branch workspaces are separately reported. Python, BLAS/eigensolver
internals and evaluator/archive storage are excluded from the array ledger.

Numerical differentiation checks the model gradient to maximum absolute
error $3.52\times10^{-10}$ and a zero-output birth preserves predictions to
$3.47\times10^{-17}$. Raw archives contain observations, teacher/evaluator
information in separate fields, all parameters and optimizer snapshots,
source/configuration hashes and resource logs. Full result reconstruction
and external novelty review are separate from these implementation checks.
The main experiment neither establishes a new theorem nor a general neural
advantage. Its contribution to the investigation is a reproducible attempt
to learn structural directions, a failed attribution control, and direct
separation of population opportunity from finite-sample nonlinear learning.
